\RequirePackage{currfile}
\documentclass[12pt]{beamer}
\usepackage[utf8]{inputenc}
\usepackage[spanish]{babel}
\usepackage{standalone}
\usepackage{color}
\usepackage{siunitx}
\usepackage{hyperref}
\usepackage[outdir=./]{epstopdf}
%\hypersetup{colorlinks,linkcolor=,urlcolor=blue}
%\hypersetup{colorlinks,urlcolor=blue}
\usepackage{xcolor,soul}
\usepackage{etoolbox}
\usepackage{amsmath}
\usepackage{amsthm}
\usepackage{mathtools}
\usepackage{tcolorbox}
\usepackage{physics}
\usepackage{multicol}
\usepackage{bookmark}
\usepackage{longtable}
\usepackage{listings}
\usepackage{cancel}
\usepackage{wrapfig}
\usepackage{empheq}
\usepackage{graphicx}
\usepackage{tikz}
\usetikzlibrary{calc, patterns, matrix, backgrounds, decorations,shapes, arrows.meta}
\usepackage[autostyle,spanish=mexican]{csquotes}
\usepackage[os=win]{menukeys}
\usepackage{pifont}
\usepackage{pbox}
\usepackage{bm}
\usepackage{caption}
\captionsetup{font=scriptsize,labelfont=scriptsize}
%\usepackage[sfdefault]{roboto}  %% Option 'sfdefault' only if the base font of the document is to be sans serif

%Sección de definición de colores
\definecolor{ao}{rgb}{0.0, 0.5, 0.0}
\definecolor{bisque}{rgb}{1.0, 0.89, 0.77}
\definecolor{amber}{rgb}{1.0, 0.75, 0.0}
\definecolor{armygreen}{rgb}{0.29, 0.33, 0.13}
\definecolor{alizarin}{rgb}{0.82, 0.1, 0.26}
\definecolor{cadetblue}{rgb}{0.37, 0.62, 0.63}
\definecolor{deepblue}{rgb}{0,0,0.5}
\definecolor{brown}{rgb}{0.59, 0.29, 0.0}
\definecolor{OliveGreen}{rgb}{0,0.25,0}
\definecolor{mycolor}{rgb}{0.122, 0.435, 0.698}

\newcommand*{\boxcolor}{orange}
\makeatletter
\newcommand{\boxedcolor}[1]{\textcolor{\boxcolor}{%
\tikz[baseline={([yshift=-1ex]current bounding box.center)}] \node [rectangle, minimum width=1ex, thick, rounded corners,draw] {\normalcolor\m@th$\displaystyle#1$};}}
 \makeatother

\newtcbox{\mybox}{on line,
  colframe=mycolor,colback=mycolor!10!white,
  boxrule=0.5pt,arc=4pt,boxsep=0pt,left=6pt,right=6pt,top=6pt,bottom=6pt}

\usefonttheme[onlymath]{serif}
%Sección de definición de nuevos comandos

\newcommand*{\TitleParbox}[1]{\parbox[c]{1.75cm}{\raggedright #1}}%
\newcommand{\python}{\texttt{python}}
\newcommand{\textoazul}[1]{\textcolor{blue}{#1}}
\newcommand{\azulfuerte}[1]{\textcolor{blue}{\textbf{#1}}}
\newcommand{\funcionazul}[1]{\textcolor{blue}{\textbf{\texttt{#1}}}}
\newcommand{\ptilde}[1]{\ensuremath{{#1}^{\prime}}}
\newcommand{\stilde}[1]{\ensuremath{{#1}^{\prime \prime}}}
\newcommand{\ttilde}[1]{\ensuremath{{#1}^{\prime \prime \prime}}}
\newcommand{\ntilde}[2]{\ensuremath{{#1}^{(#2)}}}
\renewcommand{\arraystretch}{1.5}

\newcounter{saveenumi}
\newcommand{\seti}{\setcounter{saveenumi}{\value{enumi}}}
\newcommand{\conti}{\setcounter{enumi}{\value{saveenumi}}}
\renewcommand{\rmdefault}{cmr}% cmr = Computer Modern Roman

\linespread{1.5}

\usefonttheme{professionalfonts}
%\usefonttheme{serif}
\DeclareGraphicsExtensions{.pdf,.png,.jpg}


%Sección para el tema de beamer, con el theme, usercolortheme y sección de footers
\mode<presentation>
{
  \usetheme{Warsaw}
  
  %\useoutertheme{infolines}
  \useoutertheme{default}
  \usecolortheme{crane}
  \setbeamercovered{invisible}
  % or whatever (possibly just delete it)
  \setbeamertemplate{section in toc}[sections numbered]
  \setbeamertemplate{subsection in toc}[subsections numbered]
  \setbeamertemplate{subsection in toc}{\leavevmode\leftskip=3.2em\rlap{\hskip-2em\inserttocsectionnumber.\inserttocsubsectionnumber}\inserttocsubsection\par}
  % \setbeamercolor{section in toc}{fg=blue}
  % \setbeamercolor{subsection in toc}{fg=blue}
  % \setbeamercolor{frametitle}{fg=blue}
  \setbeamertemplate{caption}[numbered]

  \setbeamertemplate{footline}
  \beamertemplatenavigationsymbolsempty
  \setbeamertemplate{headline}{}
}

\makeatletter
\setbeamercolor{section in foot}{bg=brown!75, fg=white}
\setbeamercolor{subsection in foot}{bg=cadetblue}
\setbeamertemplate{footline}
{
  \leavevmode%
  \hbox{%
  \begin{beamercolorbox}[wd=.333333\paperwidth,ht=2.25ex,dp=1ex,center]{section in foot}%
    \usebeamerfont{section in foot} \insertsection
  \end{beamercolorbox}}%
  \begin{beamercolorbox}[wd=.333333\paperwidth,ht=2.25ex,dp=1ex,center]{subsection in foot}%
    \usebeamerfont{subsection in foot}  \insertsubsection
  \end{beamercolorbox}%
  \begin{beamercolorbox}[wd=.333333\paperwidth,ht=2.25ex,dp=1ex,right]{date in head/foot}%
    \usebeamerfont{date in head/foot} \insertshortdate{} \hspace*{2em}
    \insertframenumber{} / \inserttotalframenumber \hspace*{2ex} 
  \end{beamercolorbox}}%
  \vskip0pt%
\makeatother  

\makeatletter
\patchcmd{\beamer@sectionintoc}
  {\vfill}
  {\vskip\itemsep}
  {}
  {}
\makeatother


\title{\large{Tema 3 - Bases completas y ortogonales}}
\subtitle{Objetivos}
\author{M. en C. Gustavo Contreras Mayén}
\date{\today}
\institute{Facultad de Ciencias - UNAM}
\titlegraphic{\includegraphics[width=1.75cm]{../Imagenes/escudo-facultad-ciencias}\hspace*{4.75cm}~%
   \includegraphics[width=1.75cm]{../Imagenes/escudo-unam}
}
\setbeamertemplate{navigation symbols}{}
\begin{document}
\maketitle
\fontsize{14}{14}\selectfont
\spanishdecimal{.}
\section*{Contenido}
\frame[allowframebreaks]{\tableofcontents[currentsection, hideallsubsections]}
\section{Introducción}
\frame{\tableofcontents[currentsection, hideothersubsections]}
\subsection{Objetivo del Tema}
\begin{frame}
\frametitle{Lo que hemos trabajado}
Hasta ahora hemos discutido de manera general el tipo de ecuaciones que aparecen en la física y algunas de sus propiedades, hemos trabajado con algunas de las ecuaciones diferenciales en el que estamos interesados y hemos estudiado sus singularidades.
\end{frame}
\begin{frame}
\frametitle{Lo que hemos trabajado}
Todas estas ecuaciones son ecuaciones diferenciales lineales de segundo orden, las cuales como hemos visto, admiten sólo dos soluciones linealmente independientes.
\end{frame}
\begin{frame}
\frametitle{Lo que hemos trabajado}
Nuestro objetivo en el Tema 3 es estudiar algunas propiedades adicionales de estas ecuaciones, así como algunas características generales de sus soluciones. 
\end{frame}
\begin{frame}
\frametitle{Lo que hemos trabajado}
En este tema, nos enfocaremos no a resolver la ED sino a \emph{desarrollar y comprender las propiedades generales de las soluciones}.
\end{frame}
\begin{frame}
\frametitle{Hacia dónde nos dirigimos}
Algunos de los conceptos que utilizaremos en este tema, los habrás revisado en los cursos de ecuaciones diferenciales y de álgebra lineal, por lo que en caso de que necesites darle un repaso, será conveniente para que no haya alguna complicación.
\end{frame}
\section{Temas a revisar}
\frame{\tableofcontents[currentsection, hideothersubsections]}
\subsection{Subtemas de trabajo}
\begin{frame}
\frametitle{Subtemas}
\setbeamercolor{item projected}{bg=blue!70!black,fg=yellow}
\setbeamertemplate{enumerate items}[circle]
\begin{enumerate}[<+->]
\item Ecuaciones diferenciales autoadjuntas.
\item Problemas de tipo Sturm-Liouville.
\item Operadores Hermitianos.
\item Funciones propias ortogonales.
\item Ortogonalización de Gram-Schmidt.
\item Completes de una base.
\end{enumerate}
\end{frame}
\subsection{Subtemas complementarios}
\begin{frame}
\frametitle{Temas complementarios}
Considerando que se va a requerir material adicional para tener una base sólida de conocimiento, se proponen dos temas complementarios:
\setbeamercolor{item projected}{bg=blue!70!black,fg=yellow}
\setbeamertemplate{enumerate items}[circle]
\begin{enumerate}[<+->]
\item Función delta de Dirac.
\item Espacio de Hilbert.
\end{enumerate}
\end{frame}
\begin{frame}
\frametitle{Sobre los temas complementarios}
El tema de la \emph{función delta de Dirac} es importante ya que vamos a apoyarnos bastante en lo que sigue del curso, por lo que recomendamos ampliamente que lo revisen y trabajen los ejercicios opcionales para repasar.
\end{frame}
\section{Tiempos de trabajo}
\frame{\tableofcontents[currentsection, hideothersubsections]}
\subsection{Trabajo por semana}
\begin{frame}
\frametitle{Distribución de tiempos}
Para una revisión completa de los materiales de trabajo para el Tema 3, se presenta la siguiente distribución de tiempos:
\end{frame}
\begin{frame}
\frametitle{Distribución de tiempos}
\setbeamercolor{item projected}{bg=blue!70!black,fg=yellow}
\setbeamertemplate{enumerate items}[circle]
\begin{enumerate}
\item Los subtemas Ecuaciones diferenciales autoadjuntas, Problemas de tipo Sturm-Liouville y Operadores Hermitianos.
\end{enumerate}
Se deberán de revisar en la Semana 6 (del jueves 29/10) y en la Semana 7 (del martes 3 al viernes 6 de noviembre)
\end{frame}
\begin{frame}
\frametitle{Distribución de tiempos}
\setbeamercolor{item projected}{bg=blue!70!black,fg=yellow}
\setbeamertemplate{enumerate items}[circle]
\begin{enumerate}
\item Los subtemas Funciones propias ortogonales, Ortogonalización de Gram-Schmidt y Completes de una base.
\end{enumerate}
Se deberán de revisar en la Semana 8 (del lunes 9/11 al viernes 13/11)
\end{frame}
\section{Apoyo a la calificación}
\frame{\tableofcontents[currentsection, hideothersubsections]}
\subsection{Nuevos puntajes en los ejercicios}
\begin{frame}
\frametitle{Apoyo en la puntuación (1/2)}
Se continuará con los ejercicios opcionales para cada semana, con la finalidad de incentivar el trabajo, se cambia el puntaje adicional, siendo ahora que se podrá obtener hasta $2$ puntos adicionales que se acumulará en el puntaje de la semana (pasamos de $0.5$ a $2$ puntos.)
\end{frame}
\begin{frame}
\frametitle{Apoyo en la puntuación (2/2)}
Para los dos temas complementarios se dejarán una serie de ejercicios opcionales de tal manera que podrán obtener hasta $1$ punto adicional que se acumulará en el puntaje de la semana.
\end{frame}
\begin{frame}
\frametitle{Sesiones síncronas de trabajo}
Se continuará con las sesiones de trabajo los días miércoles: 4 y 11 de noviembre a las 3 pm mediante la plataforma Zoom, considerando que ya hayan leído y revisado los materiales.
\end{frame}
\section{Punto importante}
\frame{\tableofcontents[currentsection, hideothersubsections]}
\subsection{Mensaje para considerar}
\begin{frame}
\frametitle{Avance en el curso}
Ya hemos concluido dos de los seis temas del curso, recordando que la primera evaluación corresponde a los tres primeros temas.
\end{frame}
\begin{frame}
\frametitle{Examen-Tarea}
Se hará el envío por correo con el archivo pdf con las preguntas del examen tarea correspondiente a los Temas 1 y 2.
\\
\bigskip
\pause
Con la finalidad de que los atiendan oportunamente, los resuelvan y \textbf{\textcolor{red}{los envíen de regreso en dos semanas}}.
\end{frame}
\begin{frame}
\frametitle{Acuse de recibido}
Se solicitará el acuse de recibido, es decir, esperamos una respuesta por parte de ustedes, de haber recibido el mensaje y el archivo adjunto.
\\
\bigskip
\pause
Esto nos dará la evidencia de que recibieron lo necesario para responder parte del primer examen parcial.
\end{frame}
\begin{frame}
\frametitle{Preguntas del Tema 3}
Conforme vayamos avanzando en el Tema 3, se les proporcionarán los ejercicios del examen-tarea de este tema.
\\
\bigskip
\pause
El punto es que deben de medir bien su carga de trabajo y aplicar un buen esfuerzo para estas actividades necesarias de evaluación.
\end{frame}
\end{document}